\section{Further consequences of the joint budget}\label{sec:joint-consequences}
\subsection{Separate endpoint recovery and strict gauge representations}

The signed mixture in Theorem~\ref{mol:prop:joint-budget} can be
replaced by the recovery of one conditional endpoint.  This pays for
the fact that a conditional absolute mass can be twice the mixed mass.
The next statement keeps that loss explicit.

\begin{proposition}\label{mol:prop:conditional-joint}
Under the hypotheses of Theorem~\ref{mol:prop:joint-budget}, assume
in addition that $P_Hy=0$ and $r=\|x\|>0$, and put $\Delta=R^2-r^2$.
Fix a head norming test $p=P_Hp$ with $\rho(p)=1$ and $x(p)=r$.
For every finite ancestral $K\supset H$ and each sign
$\tau\in\{-1,1\}$, choose any finite law of represented vectors
\[
 A=\sum_i a_i d_{S_i},\qquad \sum_i a_i^2\le R^2,
 \qquad \mathbb E_\tau A=P_K(x+\tau y),
\]
where each represented family consists of disjoint segments. Put
\[
 b_i=r p(S_i),\qquad A^\circ=\sum_i b_i d_{S_i},\qquad
 V=\sum_i(a_i-b_i)^2,\qquad L=r^2-\sum_i b_i^2.
\]
These laws exist by the finite atomic-ball formula. For every such
choice, the replacement satisfies
\begin{equation}\label{mol:eq:conditional-joint}
 \mathbb E_\tau(V+L)\le\Delta,
 \qquad
 \mathbb E_\tau\|Q_H(A-A^\circ)\|\le\sqrt\Delta.
\end{equation}
At each gate $g$, let $b(g)$ be the replaced crossing coefficient,
or zero when there is no crossing segment. The mixed absolute crossing
vector
$m(g)=\bigl(\mathbb E_1|b(g)|+\mathbb E_{-1}|b(g)|\bigr)/2$
satisfies $\|m\|_2\le\sqrt{2\Delta}$. Consequently a separate
endpoint argument, followed by ancestral exhaustion, yields
\begin{equation}\label{mol:eq:conditional-joint-tail}
 \|y\|\le(1+2\sqrt2)\sqrt\Delta.
\end{equation}
If the same two nonnegative errors are bounded separately, without
retaining their sum, the corresponding calculation gives
\begin{equation}\label{mol:eq:strict-gauge-tail}
 \|y\|\le(1+2\sqrt3)\sqrt\Delta.
\end{equation}
\end{proposition}

\begin{proof}
The head support of $p$ and $P_Hy=0$ imply
$P_K(x+\tau y)(p)=r$ for both signs. Expanding exactly as in
\eqref{mol:eq:joint-identity} gives the first assertion of
\eqref{mol:eq:conditional-joint}, and Cauchy--Schwarz gives the second.
For the equal mixture, put $\eta=\mathbb EV$ and
$\theta=\mathbb EL$. Then $\eta+\theta\le\Delta$, so
$\eta+2\theta\le2\Delta$.  Thus
\eqref{mol:eq:joint-absolute-mass} bounds $\|m\|_2$ as asserted.
Writing $m_\tau(g)=\mathbb E_\tau|b(g)|$, positivity gives
$m_\tau\le2m$ coordinatewise.  Hence
\[
 \|Q_H\mathbb E_\tau A^\circ\|
 \le\|m_\tau\|_2\le2\sqrt{2\Delta}.
\]
The original conditional tail is $\tau P_Ky$. Adding its replacement
error gives $\|P_Ky\|\le(1+2\sqrt2)\sqrt\Delta$ for every
finite ancestral $K\supset H$. Along an ancestral exhaustion,
$P_Ky\to y$ in norm, proving \eqref{mol:eq:conditional-joint-tail}
for the completed vector.

Alternatively, use only $\eta\le\Delta$ and $\theta\le\Delta$.
Then $\eta+2\theta\le3\Delta$, and the identical endpoint recovery
proves \eqref{mol:eq:strict-gauge-tail}.  This calculation is useful
when starting from the gauge definition. First let $x$ and $y$ have
finite coordinate support. For every
$\lambda>\max\{\|x+y\|,\|x-y\|\}$ there are finite convex
representations with coefficient radii at most $\lambda$.  Carrying
out the calculation with $\lambda$ gives
\[
 \|y\|\le(1+2\sqrt3)\sqrt{\lambda^2-r^2}.
\]
Letting $\lambda\downarrow R$ gives
\eqref{mol:eq:strict-gauge-tail}, without assuming that the infimum
in the gauge definition is attained. For a completed tail $y$, first
apply this finite-coordinate argument to $P_Ky$, where $K\supset H$
is finite and ancestral. Contractivity bounds both projected endpoint
norms by $R$, and $P_Kx=x$. The resulting estimate for $P_Ky$ passes
to $y$ along an ancestral exhaustion, as in the preceding proof.
\end{proof}

\subsection{A radial deficit budget}

A different normalization puts both endpoints in the unit ball and
replaces their coefficients by the norm-one head-test sums.  In the
notation of the original cancellation proof, let
$d=1-\|x\|$.  The joint estimate is
$\eta+\theta\le2d$, with
$\eta=\mathbb E\sum_i(a_i-p(S_i))^2$ and
$\theta=\mathbb E(1-\sum_i p(S_i)^2)$.  Consequently
\begin{equation}\label{mol:eq:radial-budget}
 \|m\|_2^2\le\eta+2\theta\le4d,
 \qquad
 \|Q_Hy\|\le(2+\sqrt2)\sqrt d.
\end{equation}
The second inequality follows by adding the mixed replacement error
$\sqrt\eta\le\sqrt{2d}$ to $\|m\|_2\le2\sqrt d$.
This retains the radial-budget route and its constants; it uses
neither individual conditional error bounds nor equality of the two
endpoint head projections.  After scaling, it reads
\[
 \|Q_Hy\|\le(2+\sqrt2)\sqrt{R(R-\|x\|)}.
\]
The coefficient $2+\sqrt2<4$ gives the convenient consequence
$\|y\|^2\le16R(R-\|x\|)$ for a zero-head displacement.
